A recursive outline tree is useful only if it carries both machine-readable structure and human-readable intent. In the Codynamic Book Machine, those two layers are kept together: the structural layer records how a node depends on other nodes, while the narrative layer explains why that dependency exists and how it should be interpreted during drafting, review, and revision.

At the machine level, a dependency is an explicit edge in the work object. A chapter, section, or subsection can declare that it builds on another node, requires another node to be complete, or should be treated as a prerequisite for downstream work. Because the same structural element model is used throughout the outline, these edges can be attached uniformly at any depth. The result is a dependency graph that is easy to inspect, compare, and route through agent workflows. A section agent can see not only where a node sits in the tree, but also which earlier material it relies on and which later material may be affected if it changes.

This structural representation matters because it separates ordering from meaning. A node appearing earlier in the outline is not necessarily a dependency, and a dependency is not merely a matter of sequence. The explicit edge says that the later node assumes concepts, definitions, or claims established elsewhere. That assumption can be checked by the gardener, used by the hypervisor when scheduling work, and surfaced to a section agent when it needs to decide whether a draft is complete enough to stand on its own.

The narrative layer complements the edge list with prose notes. These notes explain the dependency in ordinary language: what the current section inherits, what it should not repeat, and what conceptual bridge it must provide for the reader. For example, a subsection may depend structurally on a prior discussion of recursive outline trees, while its narrative note clarifies that it should focus on dependency edges rather than re-explaining the tree model itself. In this way, the dependency note acts as a local editorial contract between neighboring sections.

Narrative dependency notes are especially useful when the same structural relation has different rhetorical consequences. One section may depend on another for terminology, another for a workflow example, and another for a design principle. The machine-readable edge records the relation, but the note tells the authoring agent how to write with that relation in mind. This helps prevent repetition, reduces accidental drift, and makes it easier to preserve continuity across section boundaries.

In practice, the two forms of dependency work together. The structural edge supports automated reasoning about prerequisites, build order, and revision impact. The narrative note supports human judgment about scope, emphasis, and transition. Together they let the outline function as more than a table of contents: it becomes an executable map of the manuscript, with explicit relationships that can be inspected by agents and read by authors.

For the Codynamic Book Machine, this dual representation is part of the larger design philosophy. The system does not rely on hidden context or implicit memory to keep sections aligned. Instead, it records dependencies as durable data and explains them in prose where they matter. That makes the outline both operational and legible, which is essential when multiple agents are drafting, revising, and maintaining the same book over time.

The section is intentionally self-contained and does not require a visual to carry its argument, so no diagram is requested here.
